%% LyX 2.4.4 created this file.  For more info, see https://www.lyx.org/.
%% Do not edit unless you really know what you are doing.
\documentclass[english]{article}
\usepackage[T1]{fontenc}
\usepackage[latin9]{inputenc}
\usepackage{units}
\usepackage{enumitem}
\usepackage{amsmath}
\usepackage{amsthm}
\usepackage{geometry}
\geometry{verbose,tmargin=1in,bmargin=1in,lmargin=1in,rmargin=1in}
\PassOptionsToPackage{normalem}{ulem}
\usepackage{ulem}

\makeatletter
%%%%%%%%%%%%%%%%%%%%%%%%%%%%%% Textclass specific LaTeX commands.
\numberwithin{equation}{section}
\numberwithin{figure}{section}
\newlength{\lyxlabelwidth}      % auxiliary length 

%%%%%%%%%%%%%%%%%%%%%%%%%%%%%% User specified LaTeX commands.
\usepackage{fancyhdr}
\usepackage{lastpage}
\pagestyle{fancy}
\usepackage{enumitem}
\usepackage{ifthen}
\everymath=\expandafter{\the\everymath\displaystyle}
%\DeclareMathSizes{10}{10}{10}{10}
\usepackage{multicol}

\usepackage{tikz}
\usepackage{tikz-3dplot}
\usetikzlibrary{calc,arrows}

\usetikzlibrary{patterns}
\usetikzlibrary{plotmarks}
\usepackage{pgfplots}
\pgfplotsset{compat=newest}

\usepackage{calc}
\usepackage[nomessages]{fp}

\usepackage{datetime}
% Added by lyx2lyx
\setlength{\parskip}{\smallskipamount}
\setlength{\parindent}{0pt}

\makeatother

\usepackage{babel}
\begin{document}
\renewcommand{\author}{Dr.\,James Ashe}

\newcommand{\classNum}{331}

\newcommand{\classSec}{A}

\newcommand{\examType}{Exam 1}

\renewcommand{\date}{\formatdate{11}{9}{2013}}

%%First page's header%%%%%%%%%%%%%%%%%%%%%%
\fancypagestyle{firststyle} %%header settings for first pagr
{
	\fancyhead[L]{MTH \classNum{}(\classSec{}) \author{}\\
	\textbf{\examType{}} ~\date{}}
	\fancyhead[C]{\textbf{Name:}}
	%\fancyhead[R]{\textbf{Name:\hspace{5cm}}}
	\setlength{\headsep}{14pt}
}
%%%%%%%%%%%%%%%%%%%%%%%%%%%%%%%%%%%%%%%%%%

%%Next page's header%%%%%%%%%%%%%%%%%%%%%%%
\setlist{nolistsep}
\lhead{MTH \classNum{}(\classSec{})} 
\chead{\textbf{Name:}} 
\cfoot{\thepage{}~of~\pageref{LastPage}} 
\setlength{\headsep}{5pt} 
%%%%%%%%%%%%%%%%%%%%%%%%%%%%%%%%%%%%%%%%%%%

%%Various settings; see the preamble for other settings%%%%%
%\setenumerate[0]{leftmargin=12pt,itemindent=0pt,labelwidth=0pt}
\setlist{topsep=.5pt}
\renewcommand{\labelenumi}{\textbf{\arabic{enumi}.}}
\renewcommand{\labelenumii}{\textbf{\alph{enumii}.}}
\thispagestyle{firststyle}
%%%%%%%%%%%%%%%%%%%%%%%%%%%%%%%%%%%%%%%%%%

\newcommand{\xmax}{0}
\newcommand{\xmin}{-\xmax}
\newcommand{\xstep}{0}
\newcommand{\ymax}{0}
\newcommand{\ymin}{-\ymax}
\newcommand{\ystep}{0} 
\newcommand{\dspace}{\vspace{1cm}}
\newcommand{\xa}{0}
\newcommand{\xb}{0}
\newcommand{\ya}{1}
\newcommand{\yb}{1}

\newcommand{\xspace}{\vspace{1cm}}

\textbf{Justify your answers and show your work. Unless stated otherwise,
all questions worth 5pts.}
\begin{enumerate}[resume]
\item {[}2 pts each{]} Answer each question.

\begin{enumerate}[resume]
\item Write the infinite series, $2+4+6+8\cdots$, using sigma notation
(i.e. using $\Sigma$) . \xspace
\item Evaluate the series ${\displaystyle \sum_{k=2}^{n}(2k+1)}$ for $n=4$.\xspace
\item If $0<a_{n}<b_{n}$ and $\left\{ a_{n}\right\} $ diverges then what
can be said about $\left\{ b_{n}\right\} $?\xspace
\item If ${\displaystyle \lim_{n\rightarrow\infty}\left\{ a_{n}\right\} =0}$
then what does the Divergence Test say about${\displaystyle {\displaystyle \sum_{k=0}^{\infty}a_{k}}}$?\xspace
\item When approximating $\sqrt[3]{27.5}$ with a Taylor polynomial, what
function and center should be used. \xspace
\end{enumerate}
\end{enumerate}
\emph{Evaluate the following sequences. }
\begin{enumerate}[resume]
\item ${\displaystyle \lim_{n\rightarrow\infty}\left\{ \frac{n^{3}-2n^{2}}{5n^{3}+2n}\right\} }$\vfill
\item ${\displaystyle \lim_{n\rightarrow\infty}\left\{ \frac{0.3^{n}}{n^{10}}\right\} }$\vfill
\end{enumerate}
\emph{Evaluate the geometric series or state that it diverges. }
\begin{enumerate}[resume]
\item ${\displaystyle \sum_{k=0}^{\infty}5\frac{4^{k}}{3^{k}}}$\vfill
\item ${\displaystyle \sum_{k=2}^{\infty}3\left(\frac{1}{2}\right)^{k}}$\vfill\newpage
\end{enumerate}
\emph{Use the Divergence Test to determine whether the following series
diverge, or state that the test is inconclusive.}
\begin{enumerate}[resume]
\item ${\displaystyle \sum_{k=1}^{\infty}\frac{k^{2}-1}{3k^{2}+10}}$\vfill
\end{enumerate}
\emph{Using the Squeeze Theorem, find the limit of the following sequence
or determine that the limit does not exist.}
\begin{enumerate}[resume]
\item ${\displaystyle \lim_{n\rightarrow\infty}\left\{ \frac{(-1)^{n}}{n}\right\} }$.\vfill
\end{enumerate}
\emph{Use the Integral Test to determine whether the series converges,
diverges, or state why the test cannot be used. }\emph{\uline{Check
that the conditions of the test are satisfied}}\emph{.}
\begin{enumerate}[resume]
\item {[}10 pts{]} ${\displaystyle \sum_{k=0}^{\infty}\frac{1}{x^{\nicefrac{2}{3}}}}$.\vfill\xspace
\end{enumerate}
\emph{Use the function} $f(x)=\sqrt{x}$ \emph{to answer the following}.
\begin{enumerate}[resume]
\item Find the Taylor polynomial of order $n=3$ for $f(x)$ centered at
$1$.\vfill
\item Use the Taylor polynomial from above to approximate $\sqrt{1.05}$.\xspace\xspace
\item Compute the absolute error in the approximation assuming the exact
value is given by a calculator.\xspace
\end{enumerate}

\end{document}
